% BHU PYQ -- Manually transcribed & error-corrected
\documentclass[11pt,a4paper]{article}

\usepackage[a4paper,top=2.5cm,bottom=2.5cm,left=2.8cm,right=2.8cm,headheight=14pt]{geometry}
\usepackage{amsmath,amssymb}
\usepackage[T1]{fontenc}
\usepackage[utf8]{inputenc}
\usepackage{lmodern}
\usepackage{microtype}
\usepackage{fancyhdr}
\usepackage{enumitem}
\usepackage{hyperref}
\usepackage{lastpage}
\usepackage{parskip}

\hypersetup{colorlinks=true,urlcolor=black,linkcolor=black,
  pdftitle={CHB-505: Environmental & Nuclear Chemistry -- BHU PYQ},pdfauthor={Banaras Hindu University}}

\pagestyle{fancy}\fancyhf{}
\renewcommand{\headrulewidth}{0.4pt}\renewcommand{\footrulewidth}{0.4pt}
\fancyhead[L]{\small   \textit{Banaras Hindu University}}
\fancyhead[R]{\small   \textit{CHB-505: Environmental \& Nuclear Chemistry}}
\fancyfoot[C]{\small Page  \thepage\ of \pageref{LastPage}}


\newlist{parts}{enumerate}{2}
\setlist[parts,1]{label=(\alph*),leftmargin=*,itemsep=5pt,topsep=3pt}
\setlist[parts,2]{label=(\roman*),leftmargin=*,itemsep=3pt,topsep=2pt}

\newcommand{\pts}[1]{\hfill   \textit{[#1]}}


\begin{document}


\begin{center}
  {\large\bfseries BANARAS HINDU UNIVERSITY}\\[2pt]
  {\normalsize B.Sc. (Hons.) Semester V Examination, 2016-17}\\[6pt]
  \rule{0.8\linewidth}{0.5pt}\\[6pt]
  {\large\bfseries Paper No. CHB-505: Environmental Chemistry \& Nuclear Chemistry}\\[4pt]
  {\normalsize Time: 3 Hours\hfill Maximum Marks: 70}
\end{center}
\rule{\linewidth}{0.4pt}
\small



\noindent   \textit{Instructions: Answer five questions in total, including Question~1 from each section which is compulsory. Use separate answer books for Section~A and Section~B.}

\normalsize
\bigskip


\begin{center}
  {\large\bfseries SECTION A}\\[2pt]
  {\normalsize (Environmental Chemistry -- Marks: 35)}
\end{center}
\rule{0.4\linewidth}{0.2pt}
\smallskip




\noindent   \textbf{Question 1.} Answer all parts of the following: \pts{5 $ \times$ 2 = 10}

\begin{parts}
  \item Define the terms: tropopause and stratosphere.
  \item What is meant by a 'sink' in environmental chemistry? Give one example.
  \item What is organic soot? Write its environmental sources.
  \item List the major greenhouse gases and their relative contributions to global warming.
  \item What is the hydrosphere?
\end{parts}

\medskip\rule{\linewidth}{0.2pt}




\noindent   \textbf{Question 2.}

\begin{parts}
  \item What are the broad categories of water pollutants? What are the important water quality parameters and how do you estimate them? \pts{8}
  \item Discuss the statement: ``Ozone acts as a protective shield for living organisms on earth.'' Explain its depletion by chlorofluorocarbons (CFCs). \pts{6}
\end{parts}

\medskip\rule{\linewidth}{0.2pt}




\noindent   \textbf{Question 3.} Write detailed notes on any two of the following: \pts{2 $ \times$ 7 = 14}

\begin{parts}
  \item Photochemical Smog formation and its control
  \item Nitrogen cycle and phosphate cycle in nature
  \item Arsenic and Mercury poisoning and their toxicological effects
  \item Municipal wastewater treatment methods
\end{parts}


\newpage

\begin{center}
  {\large\bfseries SECTION B}\\[2pt]
  {\normalsize (Nuclear Chemistry -- Marks: 35)}
\end{center}
\rule{0.4\linewidth}{0.2pt}
\smallskip




\noindent   \textbf{Question 4.} Answer all parts of the following: \pts{5 $ \times$ 2 = 10}

\begin{parts}
  \item What is a nuclear reaction cross section ($\sigma$)? How is it different from the geometric cross section?
  \item Thermal neutrons interact with $^{235} \text{U}$ in two ways: (i) radiative capture ($ \text{n}, \gamma$) with $\sigma_c = 93\, \text{b}$ and (ii) fission ($ \text{n},  \text{f}$) with $\sigma_f = 584\, \text{b}$. Explain why they are different.
  \item What are the tracer applications of radioisotopes in chemistry?
  \item Differentiate between nuclear fission and nuclear fusion.
  \item Define half-life and decay constant of a radioactive isotope.
\end{parts}

\medskip\rule{\linewidth}{0.2pt}




\noindent   \textbf{Question 5.}

\begin{parts}
  \item Discuss the compound nucleus theory of nuclear reactions and outline its salient features. \pts{7}
  \item Describe the construction and working of the Geiger-M\"uller counter for the detection of nuclear radiations. \pts{7}
\end{parts}

\medskip\rule{\linewidth}{0.2pt}




\noindent   \textbf{Question 6.}

\begin{parts}
  \item Discuss the shell model of the nucleus. What are magic numbers and how do they explain nuclear stability? \pts{8}
  \item Discuss the principles and applications of neutron activation analysis (NAA). \pts{6}
\end{parts}

\vfill
\rule{\linewidth}{0.4pt}

\begin{center}\small
  \href{https://bhu-exam.vercel.app/}{Click here to visit our website}\\[4pt]
  \href{https://me-aryan.vercel.app/}{Click here to visit our contributor}\\[4pt]
  \href{https://quashan-me.vercel.app/}{Click here to visit our contributor}
\end{center}

\end{document}
